\section{Validation of the vehicle communities}
\label{app:communities}
% Discuss > Evidence tests (communities.md 4.5) -> Null model (100 randomised fleets by curveball trades; modularity z and p), seed stability (median ARI of 50 seeds), temporal stability (early against late dates, ARI and NMI) and view agreement (each view against the combined partition). Take the values of the accepted partition from evidence.parquet; the values of the earlier runs are of replaced partitions.
% Discuss > Resolution sweep -> Corridor containment and stability for each resolution. The partition is sensitive to the network near the 90% containment target (0.5 and 0.75); Bellville is selected under both.


\subsection{Similarity graph and partition}

Vehicle communities are detected with a multi-view similarity graph. Each vehicle is described by four views of non-negative feature weights, which are computed from the stop events and edges inside the study area on valid weekdays: the dwell hours at each non-home node, the vehicle-days on each network edge, the dwell hours at each formal node, and the vehicle-days with each pair of consecutive anchor dwells (origin-destination pairs). The features of each view are weighted with term frequency-inverse document frequency (TF-IDF)~\cite{salton1988_term_weighting}:
\begin{equation}
    \tilde{x}^{(m)}_{vf} = \ln\left(1 + x^{(m)}_{vf}\right)\left(\ln\frac{1 + n}{1 + \mathrm{df}_f} + 1\right),
\end{equation}

\noindent where \( x^{(m)}_{vf} \) is the weight of feature \( f \) for vehicle \( v \) in view \( m \), \( n \) is the number of vehicles, and \( \mathrm{df}_f \) is the number of vehicles that use feature \( f \). Each weight vector is subsequently normalised to unit length, such that a hub used by the whole fleet contributes less to the similarity of two vehicles than a stop shared by only a few. The combined similarity \( S_{uv} \in [0, 1] \) of vehicles \( u \) and \( v \) is the mean cosine similarity over the views that both vehicles have. A similarity graph joins each vehicle to its 15 most similar vehicles, and the Leiden algorithm~\cite{traag2019_louvain_to_leiden} partitions the graph by maximising the modularity with a resolution parameter \( \gamma \)~\cite{newman2004_community_structure,reichardt2006_statistical_mechanics_communities}:
\begin{equation}
    Q_\gamma = \frac{1}{2W} \sum_{u,v} \left[ A_{uv} - \gamma \frac{k_u k_v}{2W} \right] \mathbf{1}\left[c_u = c_v\right],
\end{equation}

\noindent where \( A_{uv} \) is the edge weight between vehicles \( u \) and \( v \), \( k_u \) is the strength of vehicle \( u \), \( W \) is the total edge weight, \( c_u \) is the community of \( u \), and \( \mathbf{1}[\cdot] \) is the indicator function. A larger \( \gamma \) produces smaller communities.

For each resolution \( \gamma \in \{0.25, 0.5, 0.75, 1, 1.25, 1.5, 2, 2.5, 3\} \), the Leiden algorithm is run with 50 seeds, and the seed partitions are combined into a consensus partition~\cite{lancichinetti2012_consensus_clustering}. A resolution is eligible when its median seed partition contains at least two communities of five or more vehicles, and when its consensus partition retains at least 90\% of corridor use within one community, where a corridor is an undirected pair of consecutive anchor dwells that recurs on at least three vehicle-days and two dates. The latter condition encodes the operational meaning of a community, namely vehicles that share corridors in a common operating area. Of the eligible resolutions, the most stable resolution, as measured by the mean pairwise adjusted Rand index (ARI)~\cite{hubert1985_comparing_partitions} of its seed partitions, is selected. Vehicles in communities of fewer than five vehicles are not assigned to a community. The existence of the communities is tested against a null model of randomised fleets~\cite{strona2014_curveball}, and their seed stability, temporal stability and agreement across views are assessed in this appendix.

\section{Network construction and coverage}
\label{app:network}
% Discuss > Support sweep (the support rule controls the captured events and dwell, not the represented distance) and pass-radius sweep (50, 100 and 150 m give a represented distance of 63.6%, 66.8% and 68.5%); the review queue (candidate cells, doubtful edges and formal/support conflicts; the speed rule only for edges with a median travel time of 160 s or more) and its decisions; the close-node distance (300 m, from a comparison of 300, 400 and 500 m).
% Discuss > Coverage measures -> Captured events and dwell, represented distance between stops, and the reasons for lost distance.

\subsection{Stop events and homes}

Stop events are extracted from the GPS trajectories with an anchored-dwell algorithm that tolerates isolated positioning errors. A candidate stop is anchored at the median position of its first three inlier points, and a later point is an inlier when it lies within 30~m of the anchor and the pairwise diameter of the inliers remains at or below 60~m. One isolated outlier is permitted when the subsequent point returns to the candidate, whereas a second consecutive outlier, or a reporting gap of more than 300~s, ends the candidate. A candidate becomes a stop event when it contains at least three inlier points, has an outlier share of at most 10\%, and lasts at least \( T_{\min} = 180 \)~s, a value that was selected from a sweep of 30 to 600~s. At the sampling interval of about 80~s, three inlier points span at least about 160~s. Consequently, stops shorter than about 2 to 3 minutes, which include most of the boarding and alighting stops along a route, cannot be observed in the data. The location of a stop event is the median of its inlier positions.

The home base of each vehicle is identified with a boundary-point rule. The first and last points of each valid vehicle-day are collected, the boundary point with the most other boundary points within \( R_{\text{home}} = 300 \)~m is selected as the candidate, and the home location is the median of the boundary points within \( R_{\text{home}} \) of the candidate. A vehicle has a home when both the first and the last points of at least 75\% of its valid vehicle-days lie within \( R_{\text{home}} \) of this location; otherwise, its most frequent first point is used as its start location. The sensitivity of the rule was assessed for radii from 150 to 1{,}000~m and minimum shares of 0.6, 0.75 and 0.9.

\subsection{Nodes, edges and review}

A single operational network is constructed for the whole study area and is used by all communities. Rather than representing the full road network, the network represents the locations at which taxis stop in an operationally important way, the fixed infrastructure that the fleet uses, and the links between these locations, such that nodes are derived from evidence of stopping and edges from evidence of driving. The eligible stop events (inside the study area, on valid weekday vehicle-days, and not at a home unless that home is at infrastructure) are counted in two square grids of 100~m cells with offsets of 0 and 50~m, so that a rank on a cell boundary is not split into two weakly supported cells. A cell is selected when it contains stop events from at least \( N_{vd} = 5 \) vehicle-days and \( N_v = 3 \) distinct vehicles. This support rule replaces the classifier of an earlier stop-network method and achieved a higher precision and F1 score than the classifier on held-out labelled cells from Gauteng. The selected cells are merged into support nodes by connected complete-linkage clustering with a maximum diameter of \( D_{\max} = 200 \)~m, which prevents long spatial chains, and each stop event is associated with one node only.

Formal nodes are obtained from the City of Cape Town minibus-taxi route termini, the taxi ranks and railway stations in OpenStreetMap (OSM), and one large rank that these sources omit, namely the Site C rank in Khayelitsha, which was confirmed on satellite imagery (Section~\ref{sec:spatial_data}). Formal nodes within 30~m of each other are merged into one node, a formal node is retained only when the fleet stops there, and a support node within the snap distance \( S = 100 \)~m of a formal node joins it and adopts its location. Home bases that are not at infrastructure become home nodes, which vehicles can reach but cannot traverse.

Since a moving vehicle covers a median distance of 435~m between consecutive GPS points, a vehicle can pass through several nodes between two stop events without a detected stop. A pass is therefore recorded when the straight line between two consecutive GPS points comes within \( r_{\text{pass}} = 100 \)~m of a node, with the pass time and cumulative distance interpolated at the point of closest approach. Passes split the movement between two stop events into hops, but they do not count towards the support of a node, so that nodes remain defined by where taxis stop. The network contains three kinds of directed edges: observed edges, which join the two nodes of a hop and are retained when at least two vehicle-days use them; derived reverse edges for observed edges without an observed counterpart; and close-node links in both directions between nodes within \( D_c = 300 \)~m of each other, which allow short movements between nearby nodes without long detours. All edge lengths are shortest road paths on the drivable OSM network, such that one-way streets are respected.

Rather than relying solely on automatic rules, the network is subjected to a single manual review of 211 items, which comprise candidate cells with high vehicle-day support but few vehicles, edges with implausible ratios of road length to straight-line length or implausible speeds, and conflicts between support nodes and formal nodes. Each item was reviewed in an interactive map that plays back the relevant vehicle-days, and 2 of 30 candidate cells, 153 of 168 flagged edges and all 13 conflicts were accepted. The resulting network comprises 726 nodes, of which 18 are ranks, and 98.9\% of the nodes lie in the largest strongly connected component. The nodes capture 76\% of the eligible stop events and 72\% of their dwell time, and 66.5\% of the distance driven between stop events runs on network edges.

\subsection{Route library and route sets}

Since the intermediate nodes of a taxi vary from trip to trip between the same terminals, the observed trips are compressed into representative routes. For each ordered terminal pair, the distinct node sequences are processed in order of decreasing support, and a sequence joins the most similar existing variant when the Jaccard similarity of their node sets~\cite{jaccard1912_flora_distribution} is at least \( \theta = 0.5 \); otherwise, it founds a new variant. A variant qualifies as a route when its trips occur on at least three vehicle-days and two dates. The representative of a qualifying variant is its medoid, i.e., the observed trip with the highest mean similarity to the other trips of the variant, and the stops of the route are the terminals and stop nodes of the medoid, together with each pass that at least half of the variant's trips visit. This majority rule removes the noise of parallel roads while retaining the nodes at which taxis consistently search for passengers. All route nodes are possible stops, and whether a driver stops at a node is a behavioural decision in \MA{} (Section~\ref{sec:abm}).

The route library is built from the trips of the whole fleet, because a corridor is a property of the operation and more vehicles provide more support. Each route receives a reverse counterpart, which is either the most similar qualified route in the opposite direction or a derived route with the reversed node sequence. Additionally, each loop is split at its turning node (the node farthest from the origin) into an outward and a return half, so that a taxi can turn at the turning node and start a new run there. The resulting library comprises 586 qualified routes and 503 derived reverse routes.

The eligible routes of a driver are the smallest set of its routes, ordered by their share of its trip distance on qualified routes, that covers 80\% of that distance, together with the reverse counterparts of these routes and the two halves of each eligible loop. In the selected communities, routes shorter than 5~km in the network are excluded, since a route whose node sequence retains only a few nodes can be much shorter in the network than its observed trips and would cause an agent to drive too little. A vehicle with at least one eligible route becomes an agent of both models. The routes carry no time of day, and the direction and timing of demand are instead handled by the demand model (Section~\ref{sec:demand}). Finally, a route is treated as a single unit of operation in \MA: once a driver has started a route, it cannot charge at a rank along that route. Consequently, an EV does not select a route that it cannot complete from a full battery (Section~\ref{sec:abm}).

\section{Agent-based model specification}
\label{app:abm}
% Discuss > ODD description of parasim-engine as configured (docs/model-agent-based.md): entities, process overview, route choice, stopping and boarding, demand, charging, memory, stochasticity and observation.
% Discuss > Engine corrections made in the study (a connector only when the policy says charge and the shift has not ended; the energy of the last charging step counts; the route energy test and the charge before going home; charging continues below the need of the cheapest route; no route that a full battery cannot finish): what each corrects, and that the ICE logs stay identical.

\subsection{Population and shifts}

The drivers of a community comprise its agents and their synthetic ICE copies. Each agent keeps its home node, its eligible routes and its idle locations, which are the terminals of its routes. Since the observed vehicles are only a sample of the taxis that operate from a hub, each agent receives \( \lfloor k - 1 \rfloor \) or \( \lceil k - 1 \rceil \) ICE copies from a seeded Bernoulli draw, such that the population is on average \( k \) times the number of agents, where \( k \) is the scale factor of Section~\ref{sec:demand}. A copy adopts the routes, idle locations and shift pattern of its agent, its home is drawn from the five network nodes nearest to the agent's home, and the start and end of its shift are moved by a common normal draw with a standard deviation of 30~min. This construction rests on the assumption that the observed agents are representative of the taxis that load at the hub of the community, both in the routes that they drive and in their working hours. The resulting populations comprise 1{,}358, 1{,}673 and 686 drivers in Khayelitsha, Bellville and Wynberg, respectively, of which the electrified agents form only 2.6\% to 6.6\% (Table~\ref{tab:models}).

Each agent has a typical-weekday shift, of which the start and end times are normally distributed with the mean and standard deviation of the first and last movements on its calendar days in the calibration period, and each day is a working day with the observed probability of activity. Two calibrated parameters, the shift spread \( \psi \) and the shift end offset \( \delta \), scale the standard deviations and move the mean end time, respectively. The offset is necessary because the observed end of a shift is the arrival at home, whereas the simulated end is the time at which the driver starts to return home. In Khayelitsha, a last-run rule additionally prevents a driver from starting a route that it expects to end after its shift end, because too many agents otherwise remained active late at night.

\subsection{Travel time}

The travel time of an edge \( \ell \) that a driver enters in time window \( w \) is modelled as follows:
\begin{equation}
    T_\ell = \frac{s_\ell}{80\ \text{km/h}} + X_{\ell w}, \qquad
    X_{\ell w} =
    \begin{cases}
        0, & \text{with probability } \pi_{\ell w}, \\
        c_w Y, \quad \ln Y \sim \mathcal{N}\left(\mu_{\ell w}, \sigma_{\ell w}^2\right), & \text{otherwise},
    \end{cases}
\end{equation}

\noindent where \( s_\ell \) is the road length of the edge, \( \pi_{\ell w} \) is the share of observed hops without excess time, \( \mu_{\ell w} \) and \( \sigma_{\ell w} \) are the parameters of a log-normal fit of the observed excess times~\cite{fosgerau2012_travel_time_variability}, and \( c_w \) is an optional excess factor. The four time windows are the night (19:00 to 06:00), the AM peak (06:00 to 09:00), the midday (09:00 to 15:00) and the PM peak (15:00 to 19:00). The fit uses the hops of the whole fleet on the calibration days, excluding hops slower than 5~km/h, since such hops contain a stand rather than driving and would make the simulated vehicles much slower than the observed vehicles. Edges with fewer than five retained hops receive a fixed excess from the pooled speed of the window, reverse edges copy the model of their observed counterpart, and close-node links use the median speed of the short observed edges. In Khayelitsha, where the simulated agents drove faster than the observed vehicles in every window, the excess is scaled by a factor \( c_w \) for each window to match the observed median speed. Drivers travel along shortest paths by road length~\cite{dijkstra1959_note_on_two_problems}, and the model does not represent congestion. Only the 18 formal ranks of the network act as engine ranks with loading queues.

\subsection{Route choice}

A driver that is not on a route first returns home when its shift has ended. Otherwise, an EV driver consults its charging policy, after which the driver chooses between its routes and idling with a multinomial logit model~\cite{mcfadden1974_conditional_logit,train2009_discrete_choice}:
\begin{equation}
    \Pr(j) = \frac{\exp\left(U_j / \tau(t)\right)}{\sum_{i} \exp\left(U_i / \tau(t)\right)},
\end{equation}

\noindent where \( U_j \) is the utility of route \( j \) or of idling, and \( \tau(t) \) is a temperature that decreases from 1 to 0.1 during an initial exploration period of ten days, such that drivers initially sample all of their routes~\cite{sutton2018_reinforcement_learning}. The utility of a route is a weighted sum of normalised features: the waiting requests at the current node that the route can carry, the competing drivers at the node on the same route, and the remembered distance, passengers and passenger distance per run of the route in the current interval of the day. The utility of idling comprises a bias \( \beta_0 \), a demand term with weight \( \beta_h \) that increases as the driver's remembered demand in the current interval decreases, and a small income term, such that drivers idle more during intervals that they remember as quiet. The route weights are fixed at the values of an earlier calibration of the model, whereas \( \beta_0 \) and \( \beta_h \) are calibrated (Section~\ref{sec:calibration}).

\subsection{Boarding and memory}

A passenger boards only a taxi that is at the origin node, has a free seat and serves a route on which the destination follows the origin; at a rank, the first compatible taxi in the loading queue is selected. Otherwise, the request waits and expires after 2~h. Passengers do not walk to other nodes or switch to other modes. The memory of each driver records the distance, duration, passengers and passenger distance of each completed route run, as well as the queue length, waiting time and charging time of each charging visit, as streaming statistics in 30-minute intervals of the day. A driver therefore learns where and when its routes carry passengers.

\subsection{Run horizon}

A cold run starts with empty memories and runs a warm-up period followed by 14 evaluation days~\cite{welch1983_statistical_analysis_simulation,law2015_simulation_modeling}. The warm-up length of each community, between 49 and 91 days, is the first day after which the seven-day means of the served passengers, their share of the requests, the occupancy and the route-choice shares remain within a tolerance of their long-run means on three seeds, rounded up to whole weeks. To reduce the computational cost, the warm-up is run once, and the resulting driver memories are saved as a warm start. Each subsequent run loads these memories, disables exploration, and runs 14 settling days followed by 14 evaluation days.

\subsection{Charging and reserve rules}

To prevent avoidable stranding, an EV maintains a reserve of \( 0.2\,C \). It does not select a route that a full battery cannot complete together with the drive from the last stop to the nearest station and the reserve. Moreover, it starts a route only when its energy covers the drive to the first stop, the route, the drive to the nearest station and the reserve, and otherwise charges first, and it charges before returning home when it would otherwise arrive below the reserve. At a station, a driver takes a connector only while it still needs to charge: before its shift end, while it is below its charge target or the energy needed for its cheapest route, or while its policy requests charging; and after its shift end, only while it would otherwise arrive home below the reserve. Home charging at 3.7~kW is applied for the time at home. An EV that nevertheless strands loses the passengers on board and is restored at home with a full battery at the next day start. These rules are corrections of the engine that replace its earlier charging behaviour without a setting to revert to it. Since ICE drivers never charge, the corrections do not change the ICE runs, which was verified by reproducing the logs of a frozen ICE run exactly.

\subsection{Passenger demand}

A departure is an anchor-to-anchor trip of an agent on the calibration days that is at least 1~km long and leaves the catchment of its origin. The catchment is defined as all nodes within 300~m of the same facility of the Comprehensive Integrated Transport Plan (CITP) of the City of Cape Town or, if no such facility exists, the origin node alone, which removes short movements within a rank complex. Since the stop rates, slow-moving shares and speeds of the observed trips do not distinguish loaded from empty legs at this sampling interval, the direction of demand is taken from survey data. Each departure receives a load factor \( \lambda \), which equals 1 except for departures in the counter-peak direction of a corridor during the AM peak (06:00 to 09:00) and in the AM-peak direction during the PM peak (15:00 to 19:00). For these departures, \( \lambda \) equals the ratio of counter-peak to peak vehicle trips on the corridor in CITP Table 3-18, which ranges between 0.09 and 0.39, or the mean ratio of 0.21 for corridors without a survey. The peak direction of a corridor is the surveyed direction with more AM-peak trips or, for corridors without a survey, the direction from its residential end, which is identified from OSM land use. Since the agents drive only their eligible routes in \MA, the departures of each agent are moreover redistributed onto its eligible routes in proportion to their load factors with a weight \( \omega_d \), such that each agent retains its total load and the supply and demand remain consistent with what the agents drive.

The scale factor is the ratio of the full population of taxis that load at the main hub to the agents:
\begin{equation}
    k = \frac{B}{Q\,L}, \qquad L = f \sum_{d \in \mathcal{H}} \lambda_d\, \omega_d, \qquad f = \frac{n_{\text{active}}}{n_{\text{valid}}\,\lvert \mathcal{D}_c \rvert},
\end{equation}

\noindent where \( B \) is the daily number of boardings at the CITP facility of the main hub (Section~\ref{sec:spatial_data}), \( \mathcal{H} \) is the set of agent departures from the hub catchment on valid vehicle-days, and \( L \) is the number of full loads of the agents per day. The day factor \( f \) converts counts on valid vehicle-days to counts per calibration day, where \( n_{\text{active}} \) and \( n_{\text{valid}} \) are the active and valid agent vehicle-days and \( \lvert \mathcal{D}_c \rvert = 6 \) is the number of calibration days, which corrects for the strong variation of the valid share between dates. This formulation assumes that the routes of the community represent all routes of its hub and that departures in the peak direction and off-peak depart full. The resulting scale factors are given in Table~\ref{tab:models}. As a plausibility check, \( k \) is compared with an upper bound that is derived from the operating licences in the provincial route register. The Bellville value exceeds this bound, which is discussed as a limitation in Section~\ref{sec:discussion_limitations}.

Each departure with load factor \( \lambda \) on a route with node sequence \( s_0, s_1, \ldots, s_m \) generates \( Q\lambda \) terminal boardings at \( s_0 \), which alight at the downstream nodes in proportion to the node weights \( w_n = 1 + c_n \), where \( c_n \) is the number of agent stop events at node \( n \) on the calibration days. The floor of one gives every route node some alightings, since short stops cannot be observed in the data. A further \( Q\lambda(1 - \alpha)/\alpha \) passengers board and alight between the en-route nodes in proportion to \( w_{s_i} w_{s_j} \), such that a share \( \alpha = 0.8 \) of all passengers boards at the terminal. The terminal share \( \alpha \) cannot be identified from the data and is therefore a judgement value. The daily number of requests of each origin-destination pair is given by:
\begin{equation}
    R_{od} = k\, f \sum_{d} \omega_d\, r_{od}(d),
\end{equation}

\noindent where the sum covers the eligible departures \( d \) on valid vehicle-days and \( r_{od}(d) \) is the number of requests that departure \( d \) generates on the pair. Since \( k \) is derived from the same departures, load factors and weights, the terminal requests at the hub equal the facility boardings by construction. Moreover, demand is generated only for servable pairs, i.e., pairs that lie in order on an eligible route of an agent.

The daily requests of each pair are distributed over the hours of the day according to the profile of its direction class (peak, counter-peak or no direction). In Wynberg, the profile is the load of the agents' eligible departures in each hour, including the evening trips home, without which the profile would have almost no evening peak. In Khayelitsha and Bellville, the departures of the agents are almost uniform over the day, whereas their moving time shows clear morning and evening peaks, i.e., at the peaks, the taxis move longer rather than more often. The profile of these communities is therefore the observed moving time of the agents in each hour, shifted earlier by half the median trip duration to approximate the boarding time and weighted by the mean load factor of the class in the period. Wynberg retains the departure-based timing with which it was calibrated and accepted. In the engine, the number of requests of each pair in each 15-minute interval is drawn from a Poisson distribution with uniformly distributed request times, which yields a non-homogeneous Poisson process with a piecewise-constant rate~\cite{ross2014_probability_models}. The same profile applies to every simulated day. The resulting demand levels are given in Table~\ref{tab:models}, and the demand is audited against the corridor survey counts and the published boardings at facilities outside the communities.

\section{Calibration targets and search}
\label{app:calibration}
% Discuss > Target values, noise and floors for each community. Strategy files with their SHA-256 values, parameters, bounds and budgets. The traffic-light assessment and the hold-out reports.
% Figure > Hourly share of agents moving in each community: the observed share (band over the resampled calibration dates) against the simulated share of the accepted model. It shows the known gap of the hourly profile (low peaks, high midday). Source: calibration assessment outputs.

\subsection{Targets and fit}

The calibration targets are computed from the observed agents on the calibration days and comprise four families: the share of agents moving in each hour of the day; the 10th, 50th and 90th percentiles of the daily distance; the same percentiles of the daily hours moving (GPS steps at 5~km/h or faster, the same threshold that separates a stand from a hop in the travel times); and the same percentiles of the coefficient of variation of each agent's daily distance. Additionally, at least 95\% of the requests must be served (a constraint), and the median seat occupancy should be as high as possible up to 0.6 (a goal), which reflects full departures from the rank and partly empty returns. The noise \( \sigma_i \) of each target is the standard deviation of its statistic over 500 bootstrap resamples of whole calibration dates~\cite{efron1993_bootstrap}, which preserves the dependence between vehicles on the same date. The fit of target \( i \) is given by:
\begin{equation}
    z_i = \frac{\left| y_i^{\text{sim}} - y_i^{\text{obs}} \right|}{\max\left(\sigma_i,\ \epsilon_i\right)},
\end{equation}

\noindent where \( y_i^{\text{sim}} \) and \( y_i^{\text{obs}} \) are the simulated and observed values of the target, and \( \epsilon_i \) is a floor (0.01 for the movement share, 5~km, 0.02 for the coefficient of variation and 0.25~h) that prevents targets with very small noise from dominating the fit. The loss of a family is the mean fit of its targets, and the loss of a run is the weighted mean of the family losses.

\subsection{Search}

Following~\cite{thiele2014_abm_parameter_estimation}, the calibration proceeds in two stages. First, hand tuning from the warm start locates a starting point within the region of plausible operation and identifies gaps that no behavioural parameter can close. Such gaps required changes to model rules, namely the travel-time scaling and the last-run rule in Khayelitsha. Second, a search strategy is frozen in a versioned file whose SHA-256 digest is recorded. The strategy specifies nine parameters and their bounds (the idle bias and idle demand weight; the stop bias and maximum stop duration; the first-stop bias, demand weight and maximum duration; and the shift spread and shift end offset), the objectives, the constraint and the budget. The Khayelitsha strategy weights the movement family by 0.25 relative to the duty-cycle families, since a search with equal weights traded the duty cycles for the hourly shape of the day, whereas the Bellville and Wynberg strategies use equal weights and apply this priority at the acceptance. The search uses the multivariate tree-structured Parzen estimator (TPE) of Optuna~\cite{bergstra2011_hyperparameter_optimisation,akiba2019_optuna,ozaki2022_multiobjective_tpe} with two objectives, namely the minimisation of the loss and the maximisation of the occupancy capped at 0.6, subject to the served-share constraint. Each trial runs one seed from the warm start, and the search starts with a batch of 40 trials, of which 10 are space-filling start trials, followed by optional continuation batches of 20 trials. Selected points on the feasible front are subsequently validated with cold runs, which include the full warm-up, on three fresh seeds.

\subsection{Assessment and hold-out}

Each validated point is assessed with graded levels (good, acceptable or limited), which are expressed relative to the loss of the observed agents against their own targets over resampled calibration dates, so that the losses of the three communities are comparable. The acceptance prioritises realistic behaviour and the duty-cycle families, followed by the stability across seeds, the served share and the occupancy, and finally the hourly shape of the day. After the acceptance, a hold-out report compares the fit of the accepted model against the targets of the hold-out days (2 and 10 December) with its fit against the calibration targets. The resulting hold-out gap is assessed against the loss of the observed hold-out days against the calibration targets, which quantifies how much two days differ from six days in the data alone. The hold-out days test only the calibrated behaviour, since the network, the communities and the routes use all study weekdays.

% TODO: The Bluedot replication is not part of the new pipeline (the study uses only the Sebenza Cape Town fleet). Rerun it with the same steps, or remove this appendix.
\section{Replication with a second fleet}
% Discuss > Bluedot, November 2022 (28 days, 621 vehicles, 05:00--20:00 only): 23 communities, seed ARI 1.0, early and late ARI 0.92; screening classes. Night is not observed, so overnight-only viability is optimistic.

\section{Reproducibility}
\label{app:reproducibility}
% Discuss > The 13 steps, their manifests (input, settings, code and output digests) and the human decision of each step, mapped to the sections of this paper. The parasim-engine commit and the SHA-256 of the engine executable in the experiments inputs. No paths, pseudonyms or coordinates.
